% GENERATED FILE. Edit canonical lessons, then run npm run generate:books.
% canonical-source-hash: fnv1a64:512402af046552a1
% canonical-lessons: FR-C449-fleuriste, FR-C449-buraliste, FR-C449-affranchir, FR-C449-expedier, FR-C449-ficeler

\chapter{Fleurs, tabac et colis}
\label{ch:fr-fleurs-tabac-et-colis}
\hlchaptermodality{449}

\begin{chapteropening}
\textbf{By the end of this chapter:} \emph{I can use local shops and send things: the florist, the tobacconist, putting a stamp on, sending off and tying up a parcel.}

Complete the last lesson of chapter 449: I can use local shops and send things: the florist, the tobacconist, putting a stamp on, sending off and tying up a parcel.
\end{chapteropening}

\section[le fleuriste]{le fleuriste — the florist}

\label{lesson:FR-C449-fleuriste}



\begin{quote}
Before the new one: say the French for the watchmaker, then the French for the jeweller.
\end{quote}

\subsection*{You'll want to know: le fleuriste}
\textbf{le fleuriste} — \textquotedblleft{}the florist\textquotedblright{}.

\begin{quote}
Le fleuriste a fait un joli bouquet pour ma mère.
\end{quote}

\begin{quote}
\emph{The florist made a pretty bunch of flowers for my mother.}
\end{quote}

Same form for a woman: \emph{la fleuriste}.

\begin{cousinweb}
\emph{fleuriste} is built on \emph{fleur} 'flower' with \emph{-iste}; \emph{fleur} is inherited from Latin \emph{flōs, flōris} (EtymDB).
\end{cousinweb}

\begin{grammarlens}[title={What you've built}]
One more word for telling, explaining and describing at B1.
\end{grammarlens}

\subsection*{Your turn}
\begin{itemize}
\raggedright
  \item \emph{Say it:} \emph{le fleuriste}
  \item \emph{Say it:} \emph{le fleuriste}, once more
  \item \emph{Say it:} a sentence of your own with \emph{le fleuriste}
  \item \emph{Recall:} say the French for the watchmaker, then the French for the jeweller, then say \emph{le fleuriste} again
\end{itemize}

\begin{tcolorbox}[breakable,colback=teal!4,colframe=teal!35!black,title={Before you move on}]
What does le fleuriste mean? (\textquotedblleft{}The florist\textquotedblright{}.) Say it once more.
\end{tcolorbox}
\section[le buraliste]{le buraliste — the tobacconist}

\label{lesson:FR-C449-buraliste}



\begin{quote}
Before the new one: say the French for the jeweller, then the French for the florist.
\end{quote}

\subsection*{You'll want to know: le buraliste}
\textbf{le buraliste} — \textquotedblleft{}the tobacconist\textquotedblright{}.

\begin{quote}
Le buraliste vend aussi des timbres et des journaux.
\end{quote}

\begin{quote}
\emph{The tobacconist also sells stamps and newspapers.}
\end{quote}

He runs \emph{le bureau de tabac}.

\begin{cousinweb}
\emph{buraliste} is built on \emph{bureau} 'office, desk' with \emph{-iste} (EtymDB). \emph{Bureau} first meant a coarse woollen cloth that covered a desk.
\end{cousinweb}

\begin{grammarlens}[title={What you've built}]
One more word for telling, explaining and describing at B1.
\end{grammarlens}

\subsection*{Your turn}
\begin{itemize}
\raggedright
  \item \emph{Say it:} \emph{le buraliste}
  \item \emph{Say it:} \emph{le buraliste}, once more
  \item \emph{Say it:} a sentence of your own with \emph{le buraliste}
  \item \emph{Recall:} say the French for the jeweller, then the French for the florist, then say \emph{le buraliste} again
\end{itemize}

\begin{tcolorbox}[breakable,colback=teal!4,colframe=teal!35!black,title={Before you move on}]
What does le buraliste mean? (\textquotedblleft{}The tobacconist\textquotedblright{}.) Say it once more.
\end{tcolorbox}
\section[affranchir]{affranchir — to put a stamp on, to frank; to free}

\label{lesson:FR-C449-affranchir}



\begin{quote}
Before the new one: say the French for the florist, then the French for the tobacconist.
\end{quote}

\subsection*{You'll want to know: affranchir}
\textbf{affranchir} — \textquotedblleft{}to put a stamp on, to frank; to free\textquotedblright{}.

\begin{quote}
Elle a affranchi la lettre au tarif normal.
\end{quote}

\begin{quote}
\emph{She put a normal-rate stamp on the letter.}
\end{quote}

\emph{L'affranchissement} — the postage. First 'to set free'.

\begin{cousinweb}
\emph{affranchir} is built from \emph{a-} and \emph{franc} 'free', the word behind \emph{franchir} (EtymDB links the two); \emph{franc} comes from the name of the Franks: a letter with its stamp travels free.
\end{cousinweb}

\begin{grammarlens}[title={What you've built}]
One more word for telling, explaining and describing at B1.
\end{grammarlens}

\subsection*{Your turn}
\begin{itemize}
\raggedright
  \item \emph{Say it:} \emph{affranchir}
  \item \emph{Say it:} \emph{affranchir}, once more
  \item \emph{Say it:} a sentence of your own with \emph{affranchir}
  \item \emph{Recall:} say the French for the florist, then the French for the tobacconist, then say \emph{affranchir} again
\end{itemize}

\begin{tcolorbox}[breakable,colback=teal!4,colframe=teal!35!black,title={Before you move on}]
What does affranchir mean? (\textquotedblleft{}To put a stamp on, to frank; to free\textquotedblright{}.) Say it once more.
\end{tcolorbox}
\section[expédier]{expédier — to send off, to ship}

\label{lesson:FR-C449-expedier}



\begin{quote}
Before the new one: say the French for the tobacconist, then the French for to put a stamp on.
\end{quote}

\subsection*{You'll want to know: expédier}
\textbf{expédier} — \textquotedblleft{}to send off, to ship\textquotedblright{}.

\begin{quote}
Nous avons expédié le colis mardi.
\end{quote}

\begin{quote}
\emph{We sent the parcel off on Tuesday.}
\end{quote}

\emph{L'expéditeur} — the sender.

\begin{cousinweb}
\emph{expédier} is built on \emph{expédient}, from Latin \emph{expedīre} 'to free the feet, to get things ready', from \emph{pēs} 'foot' (Lewis \& Short).
\end{cousinweb}

\begin{grammarlens}[title={What you've built}]
One more word for telling, explaining and describing at B1.
\end{grammarlens}

\subsection*{Your turn}
\begin{itemize}
\raggedright
  \item \emph{Say it:} \emph{expédier}
  \item \emph{Say it:} \emph{expédier}, once more
  \item \emph{Say it:} a sentence of your own with \emph{expédier}
  \item \emph{Recall:} say the French for the tobacconist, then the French for to put a stamp on, then say \emph{expédier} again
\end{itemize}

\begin{tcolorbox}[breakable,colback=teal!4,colframe=teal!35!black,title={Before you move on}]
What does expédier mean? (\textquotedblleft{}To send off, to ship\textquotedblright{}.) Say it once more.
\end{tcolorbox}
\section[ficeler]{ficeler — to tie up (with string)}

\label{lesson:FR-C449-ficeler}



\begin{quote}
Before the new one: say the French for to put a stamp on, then the French for to send off.
\end{quote}

\subsection*{You'll want to know: ficeler}
\textbf{ficeler} — \textquotedblleft{}to tie up (with string)\textquotedblright{}.

\begin{quote}
Le vendeur a ficelé le paquet avec soin.
\end{quote}

\begin{quote}
\emph{The shop assistant tied up the parcel carefully.}
\end{quote}

\emph{La ficelle} — string; also a thin baguette.

\begin{cousinweb}
\emph{ficeler} is built on \emph{ficelle} 'string' (EtymDB lists the verb), which is usually traced to Vulgar Latin \emph{fūnicella}, a small \emph{fūnis} 'rope'.
\end{cousinweb}

\begin{grammarlens}[title={What you've built}]
One more word for telling, explaining and describing at B1.
\end{grammarlens}

\subsection*{Your turn}
\begin{itemize}
\raggedright
  \item \emph{Say it:} \emph{ficeler}
  \item \emph{Say it:} \emph{ficeler}, once more
  \item \emph{Say it:} a sentence of your own with \emph{ficeler}
  \item \emph{Recall:} say the French for to put a stamp on, then the French for to send off, then say \emph{ficeler} again
\end{itemize}

\begin{tcolorbox}[breakable,colback=teal!4,colframe=teal!35!black,title={Before you move on}]
What does ficeler mean? (\textquotedblleft{}To tie up (with string)\textquotedblright{}.) Say it once more.
\end{tcolorbox}
